\documentclass[11pt,a4paper]{article}

% Kept minimal so the CI TeX Live install always compiles this.
\usepackage[T1]{fontenc}
\usepackage[utf8]{inputenc}
\usepackage[margin=1in]{geometry}
\usepackage{amsmath}
\usepackage{amssymb}
\usepackage{hyperref}

\title{zkTGuard: Privacy-preserving Telegram credentials on TON \\
       \large via zkTLS proxy attestation and Groth16}
\author{Igor Prostov}
\date{v0.1 draft --- May 2026}

\begin{document}
\maketitle

\begin{abstract}
zkTGuard is an open-source research project that turns Telegram's
existing trust signals into privacy-preserving on-chain credentials
on TON. A proxy attestor witnesses the response to a Telegram-shaped
HTTP endpoint, signs the transcript with EdDSA over BabyJubjub, and a
Groth16 circuit over BLS12-381 attests in zero knowledge that the
transcript decrypts to a JSON object whose ``account creation''
timestamp predates a public threshold. A FunC verifier on TON checks
the proof using the chain's BLS pairing precompile and records a
per-application nullifier.
\end{abstract}

\section{Introduction}\label{sec:intro}

Telegram Mini Apps host hundreds of millions of users on TON. Sybil
resistance for these apps is an unsolved problem: existing identity
systems (Human Passport, World ID, Polygon ID) ignore the
platform-native trust signals that already exist inside Telegram.
zkTGuard is, to the authors' knowledge, the first published design
that uses Telegram itself as a credential issuer via zkTLS.

The design has four moving parts: a proxy attestor (a stateless Go
service), a prover (a stateless Rust service), a FunC verifier
contract, and a soulbound credential contract. A client-side Mini App
connects them via TON Connect 2.0. The bundle is published under MIT,
with no hosted service, no telemetry, no fee handling, and no token.

\section{Related work}\label{sec:related}

Groth16 remains the canonical choice for circuits with stable
verification keys. Universal-setup systems (PLONK, Plonky3, Halo2,
Binius) avoid per-circuit ceremonies but pay a runtime cost. v0.1
picks Groth16 because the on-chain verifier reduces to a single
\texttt{BLS\_PAIRING}.

Proxy attestation follows the Reclaim Protocol construction: a notary
forwards TLS records, signs the encrypted transcript, and the prover
later proves in zero knowledge that a witnessed key decrypts to the
required field. TLSNotary provides a stronger MPC-TLS variant we view
as the natural v1 upgrade.

Identity primitives (Semaphore, zkLogin, Sui zkLogin, World ID) all
build credentials atop a non-revealing identity primitive. zkTGuard
reuses the same idea --- nullifiers as the on-chain primary key ---
and adapts it to the Telegram + TON deployment surface.

Folding schemes (Nova, HyperNova) offer a path to session continuity
that the v0 per-action proof does not yet exploit. We treat them as
future work.

\section{Cryptographic construction}\label{sec:construction}

Let $r$ be the order of the BLS12-381 scalar field, $G$ a generator
of $\mathbb{G}_1$, $H$ a generator of $\mathbb{G}_2$, and $e:
\mathbb{G}_1 \times \mathbb{G}_2 \to \mathbb{G}_T$ a Type-3 Ate pairing.
A trusted setup produces scalars $(\alpha, \beta, \gamma, \delta) \in
(\mathbb{F}_r^\times)^4$ and a chain of $\mathrm{IC}_i \in \mathbb{G}_1$
for $i \in \{0, \dots, n\}$ where $n$ is the number of public inputs.

The on-chain verifier accepts a proof $\pi = (A, B, C) \in
\mathbb{G}_1 \times \mathbb{G}_2 \times \mathbb{G}_1$ and a public-input
vector $\mathbf{x} \in \mathbb{F}_r^n$ iff
$$
e(A, B) = e(\alpha G, \beta H) \cdot e(\mathrm{vk}_x, \gamma H) \cdot e(C, \delta H),
$$
where $\mathrm{vk}_x = \mathrm{IC}_0 + \sum_{i=1}^{n} x_i \cdot \mathrm{IC}_i$.

The FunC verifier reduces this to a single four-pair pairing check.

The v0.1 circuit publishes eight scalars in this order: \texttt{nonce},
\texttt{app\_id}, \texttt{expiration}, \texttt{claim\_type},
\texttt{nullifier}, \texttt{attestor\_pubkey\_x},
\texttt{attestor\_pubkey\_y}, and \texttt{threshold\_months}. The
high 32 bits of \texttt{nonce} carry the \texttt{chain\_id} so the
verifier can enforce replay protection without a separate public
input.

The nullifier is defined as $\mathrm{nullifier} = \mathrm{Poseidon}_3(
\mathit{user\_secret}, \mathit{app\_id}, \mathit{claim\_type})$.
Per-application, per-claim, deterministic, unlinkable across
\texttt{app\_id} for the same user.

The construction follows the proxy-attestation model. The attestor
must not collude with the prover. v0.1 demonstrates with a single
test attestor; the architecture admits a future attestor set with
diverse operators using a threshold or set-membership check on chain.

\section{Implementation}\label{sec:impl}

The reference implementation is a monorepo containing a Go attestor,
a Rust prover, four FunC contracts, a Circom circuit, a TypeScript
SDK, a Mini App, and three example Mini Apps. All components are
MIT-licensed.

The attestor is a stateless Go binary that hashes the upstream
response with Poseidon over BN254 and signs with EdDSA on BabyJubjub
using \texttt{iden3/go-iden3-crypto}. It exposes \texttt{/attest},
\texttt{/pubkey}, and \texttt{/health}.

The prover is a stateless Rust binary; v0.2 supplements it with a
Node HTTP wrapper (\texttt{prover/server}) so the SDK posts witness
data to a configurable \texttt{proverUrl} and receives the
BLS12-381-compressed proof in the wire format the verifier reads.
v0.1's synthetic-proof crafter is retired; the SDK no longer ships a
local prover.

The verifier is a FunC contract that embeds the real v0.2 Phase-A VK
as compile-time constants. It reads a proof and a public-input snake
from the message body, computes the verification key linear
combination using two \texttt{BLS\_G1\_MULTIEXP} calls joined by
\texttt{BLS\_G1\_ADD}, and invokes \texttt{BLS\_PAIRING} for the
four-pair check.

\subsection{End-to-end verification on testnet}

A Playwright suite in \texttt{miniapp/e2e/} drives the Mini App
against the live Phase~D1 deployment under a sandbox wallet stub
that signs with a hardcoded testnet seed and broadcasts via
\texttt{@ton/ton}'s \texttt{TonClient}. Four scenarios run:
\texttt{happy-path} (connect, prove, mint, nullifier recorded);
\texttt{expired-proof} (verifier rejects with exit~403 before the
pairing check); \texttt{replayed-nullifier} (second submission of
the same Poseidon hash rejected with exit~402 on registry reply);
and \texttt{unregistered-app} (verifier parks, registry replies
\texttt{is\_registered=false}, parked entry dropped, no mint). The
suite is wired into a separate GitHub Actions workflow,
\texttt{e2e-testnet}, that runs on \texttt{main} and on \texttt{v*}
tag pushes; it is the gate on release tagging.

\section{Evaluation}\label{sec:eval}

Sandbox numbers from \texttt{@ton/sandbox} v0.32 on \texttt{@ton/core}
0.62 are the cheap-to-measure baseline. Testnet numbers come from a
live deployment of v0.2 Phase~D1 on TON testnet (workchain 0,
deployer wallet \texttt{kQBEhiW…XUS}). Per-hop transaction hashes are
recorded in \texttt{contracts/deployments/sanity-check.json} and
\texttt{contracts/deployments/v0.2-testnet.json}.

\begin{center}
\begin{tabular}{lrr}
Operation & Sandbox gas & Testnet \\\hline
Pairing sanity check, good pair & 62\,741 & --- \\
Pairing sanity check, bad pair  & 62\,625 & --- \\
\texttt{op::verify\_claim} accept (verify + park) & 151\,955 & confirmed in tx \texttt{3039fdf8…} \\
\texttt{op::query\_registered} (verifier $\rightarrow$ registry) & --- & confirmed in tx \texttt{9774420f…} \\
\texttt{op::query\_reply}      (registry $\rightarrow$ verifier) & --- & confirmed in tx \texttt{69558695…} \\
\texttt{op::mint}              (verifier $\rightarrow$ collection) & --- & confirmed in tx \texttt{077496a2…} \\
End-to-end (verify $\rightarrow$ mint, wall-clock) & --- & 30--60\,s under public toncenter \\
\end{tabular}
\end{center}

Per-hop gas in nano-TON on testnet was not extracted in the v0.2
sanity-check pass because the public toncenter rate limit dominated
the script's wall time; a follow-up rerun against a paid API key (or
a single \texttt{tonapi} query per transaction) will fill in the
column. Published comparators for Reclaim and zkLogin are pending a
literature search; \textbf{TODO}: cite ePrint numbers in a v0.3 pass.

\section{Limitations}\label{sec:limits}

The verifier contract embeds the real Phase-A verification key
generated by a deterministic single-party Phase-2 contribution
against a local BLS12-381 Powers of Tau. v0.1's "placeholder VK"
caveat is retired; the standing limitation is now that the trusted
setup is single-party, so a malicious operator with the same
deterministic seed can forge proofs. v0.3 will participate in (or
coordinate) a multi-party ceremony.

The attestor signs with BabyJubjub EdDSA, not BLS12-381 as the bible's
Part C originally specifies. BLS verification inside a Groth16/bn128
circuit costs millions of constraints; BabyJubjub clocks in around
four thousand.

The v0.2 release wires the cross-contract async flow: the verifier
parks an accepted proof, sends a query to the App Registry, and
completes the soulbound-credential mint on a positive reply. The
nullifier is recorded only after the reply lands, so a race-free
window for registry mutation is preserved.

The Powers-of-Tau choice is documented but no multi-party Phase-2
contribution has been performed; v0.2 ships a deterministic local
contribution. v0.2 still runs against a stub server rather than real
Telegram traffic.

\section{ChaCha20-Poly1305 in the v0.2 Phase B circuit}\label{sec:chacha}

Phase B replaces the v0.1 AEAD stubs with a real
ChaCha20-Poly1305 decryption check inside the circuit. The full
design --- quarter-round arithmetic, the five-limb Poly1305
accumulator, the AEAD composition per RFC 8439, and the
constraint-count budget --- lives in
\texttt{paper/design/chacha20-circuit.md}, kept separate from this
manuscript so it can iterate independently. The cipher-suite
constraint forces the attestor to negotiate ChaCha20-Poly1305 with
the upstream endpoint; AES-GCM TLS sessions are not provable in this
construction.

\section{Future work}\label{sec:future}

Migrate proxy attestation to MPC-TLS (TLSNotary) once that ecosystem
matures. Replace the per-action proof with folded proofs (Nova,
HyperNova) for session continuity. Standardise the claim catalog
beyond v0's \texttt{account\_age}. Compose multiple verified claims
into a higher-assurance attestation. Provide a formal security proof
of the proxy-attestation construction. Benchmark mobile-side proving
across Android and iOS for varying circuit size.

\appendix
\section{Reproducibility}\label{app:repro}

The exact source for this draft lives at
\url{https://github.com/Igorprostoff/zktguard}. To rebuild every
committed artifact from a clean clone:

\begin{verbatim}
git clone https://github.com/Igorprostoff/zktguard.git
cd zktguard
bash scripts/reproduce.sh
\end{verbatim}

The script and its expected output hashes are documented in
\texttt{REPRODUCE.md} at the repo root.

\begin{thebibliography}{99}

\bibitem{Groth16}
J. Groth.
\emph{On the size of pairing-based non-interactive arguments}.
EUROCRYPT 2016.

\bibitem{Reclaim}
Reclaim Protocol.
\emph{zkTLS via proxy attestation}.
\url{https://reclaimprotocol.org}.

\bibitem{TLSNotary}
PSE TLSNotary.
\emph{MPC-TLS for selective disclosure}.
\url{https://tlsnotary.org}.

\bibitem{zkLogin}
Mysten Labs.
\emph{zkLogin: an OAuth-friendly account scheme on Sui}.
2023.

\bibitem{SuiZkLogin}
F. Baldimtsi et al.
\emph{Sui zkLogin paper}.
IACR ePrint, 2024.

\bibitem{Semaphore}
PSE.
\emph{Semaphore: anonymous signaling on Ethereum}.
\url{https://semaphore.pse.dev}.

\bibitem{WorldID}
Worldcoin.
\emph{World ID protocol}.
\url{https://worldcoin.org/world-id}.

\bibitem{Nova}
A. Kothapalli, S. Setty, I. Tzialla.
\emph{Nova: Recursive zero-knowledge arguments from folding schemes}.
CRYPTO 2022.

\bibitem{Plonky3}
Polygon Zero.
\emph{Plonky3 specification}.
\url{https://github.com/Plonky3/Plonky3}.

\bibitem{Binius}
B. Diamond, J. Posen.
\emph{Succinct arguments over towers of binary fields}.
IACR ePrint 2023/1784.

\end{thebibliography}

\end{document}
